\documentclass[10pt,a4paper]{amsart}
\usepackage[margin=19mm]{geometry}
\usepackage{amsmath,amssymb,mathtools}
\usepackage[scr=rsfs,cal=euler]{mathalfa}
\usepackage{xfrac}
\usepackage[unicode,hidelinks,bookmarks=false]{hyperref}
\newcommand{\ie}{i.e.\ }
\newcommand{\cf}{cf.\ }
\newcommand{\resp}{respectively\ }
\newcommand{\Subsection}{\subsection}
\providecommand{\nobreakdash}{\nobreak\hbox{-}}
\setlength{\emergencystretch}{2em}
\multlinegap0pt
\usepackage{xcolor}
\usepackage{float}
\usepackage{amssymb}
\usepackage{tikz}
\usepackage{mathtools}
\usepackage{dsfont}
\newtheorem{theorem}{Theorem}
\title{Multi-graded generic initial ideals, regularity, and the optimal colorful fractional Helly theorem for $d$-Leray complexes}
\author{Daniel McGinnis}
\date{Source: arXiv:2608.25891v4 (CC BY 4.0)}
\begin{document}
\maketitle

\noindent\emph{Source and licence.} Multi-graded generic initial ideals, regularity, and the optimal colorful fractional Helly theorem for $d$-Leray complexes. Version arXiv:2608.25891v4 (CC BY 4.0), distributed by arXiv under the Creative Commons Attribution 4.0 International licence, http://creativecommons.org/licenses/by/4.0/, as recorded in the arXiv licence metadata for this exact version and shown on the abstract page https://arxiv.org/abs/2608.25891v4. Full licence notice: \texttt{licenses/arxiv-2608.25891-CC-BY-4.0.txt}.\par\smallskip

\noindent\textit{Editorial scope.} Counted unit: the article's theorem Gen (source0.tex lines 408--418). The article's complete original proof of it is delivered with it (source0.tex lines 419--454, one \texttt{proof} environment of 36 lines). No mathematics is written, repaired or re-derived: the statement and the proof are the article's own lines, in the article's own notation, and no retained line is substituted or re-flowed.  No further source-local context is needed: every cross-reference of every retained line resolves inside the two delivered spans.  Nothing was omitted from any retained span, so the package carries no omission record.  No retained line contains a citation, so the wrapper carries no bibliography and no bibliography entry.  No retained line contains a figure environment, an image-inclusion command or drawing code, so no image is delivered and the package declares no asset dependency.  Editorial formatting only, in addition: an AMS \texttt{amsart} preamble that restates the article's own declarations and macros used by the retained lines, namely the environments \texttt{theorem} on separate counters, together with the standard packages those lines need.  The retained environments print in this document as Theorem 1 in order of appearance, whereas the article prints the same items with the numbering of its own full document.  The complete native source of the article as distributed in the arXiv e-print for this exact version is bundled unchanged as \texttt{sources/208602/source0.tex}.
\begin{theorem}\label{thm:Ext Betti Gen}
    Let $\ell_1,\dots,\ell_n$ be an almost regular sequence of $M$ with degrees $k_1\leq \cdots \leq k_n$. Then
    \begin{itemize}
        \item[(a)] $m_j = \textrm{max}\{ s_1,\dots,s_j\}$ for $j=1,\dots,n$. In particular, $m_1\leq \cdots \leq m_n$.
        \item[(b)] Let $i_1<\cdots<i_p$ be the indices $i$ such that $m_i-m_{i-1} > 0$. Then for all $1\leq t \leq p$ and all $j\geq i_t$ we have:
        \begin{itemize}
            \item[(i)] $H_i(j)_{k_j+\cdots+k_{j-i+1}+s}=0$ for $s>m_{i_{t-1}}$ and $i>j-i_t+1$,
            \item[(ii)] $H_{j-i_t+1}(j)_{k_j+\cdots+ k_{i_t}+m_{i_t}} \cong (0:_{M\langle i_t -1 \rangle} \ell_{i_t})_{s_{i_t}}$.
        \end{itemize}
    \end{itemize}
\end{theorem}

\begin{proof}
     For (a), we show by induction on $j$ that $m_j=\textrm{max}\{m_{j-1},s_j\}$. For $j=1$ we have $H_i(1)=0$ for $i>1$ and $H_1(1) \cong (0:_M \ell_1)(-k_1)$, so $m_1 = s_1$ as desired.

    Let $j>1$ and assume that $s>\textrm{max}\{m_{j-1},s_j \}$. From the long exact sequence, we have the exact sequence
    \begin{equation}\label{eq:1st exact Sequence}
        H_i(j-1)_{k_j+\cdots + k_{j-i+1} + s} \longrightarrow H_i(j)_{k_j+\cdots + k_{j-i+1} + s} \longrightarrow H_{i-1}(j-1)_{k_{j-1}+\cdots + k_{j-i+1} + s}
    \end{equation}
    for all $i\geq 1$ (when $i=1$, we take $H_{i-1}(j-1) = (0:_{M\langle j-1 \rangle} \ell_j)$). Working inductively, the left and right homology groups are $0$ (for the left homology group, we are using the fact that $k_j\geq k_{j-i}$), so $H_i(j)_{k_j+\cdots + k_{j-i+1} + s}=0$ as well. 

    Now suppose that $m_j < \textrm{max}\{m_{j-1},s_j\}$. In the case that $m_j < m_{j-1}$, let $i\geq 1$ be an integer such that $H_i(j-1)_{k_{j-1}+\cdots+k_{j-i}+m_{j-1}}\neq 0$. From the long exact sequence, we have the exact sequence
    \[
    H_{i+1}(j)_{k_{j}+\cdots+k_{j-(i+1)+1}+m_{j-1}} \longrightarrow H_i(j-1)_{k_{j-1}+\cdots+k_{j-i}+m_{j-1}}\longrightarrow H_i(j-1)_{k_{j}+\cdots+k_{(j-1)-i+1}+m_{j-1}}.
    \]
    This is a contradiction since, again, the left and right homology groups are $0$.

    In the case $m_j<s_j$, the exact sequence
    \[
    H_1(j)_{k_j+s_j} \longrightarrow (0:_{M\langle j-1 \rangle} \ell_j)_{s_j} \longrightarrow 0
    \]
    implies that $(0:_{M\langle j-1 \rangle} \ell_j)_{s_j}=0$, a contradiction.

    For (b)(i), we start with the case $j=i_t$ and work inductively. We have $j-1 < i_t$, so $m_{j-1} = m_{i_{t-1}}$. Therefore, $H_i(j-1)_{k_{j-1} + \cdots k_{j-i}+s}=0$ for $s>m_{i_{t-1}}$ and $i\geq 1$ by (a). The sequence (\ref{eq:1st exact Sequence}) then implies that $H_i(j)_{k_j+\cdots k_{j-i+1} + s}=0$ for $s>m_{i_{t-1}}$ and $i>1=j-i_t+1$. When $j>i_t$, $j-1\geq i_t$, and we have by induction that $H_i(j-1)_{k_{j} + \cdots + k_{j-i+1}+s}=H_{i-1}(j-1)_{k_{j-1} + \cdots + k_{j-i+1}+s}=0$ when $s>m_{i_{t-1}}$ and $i>j-i_t+1$.

    For (b)(ii), first note that we have the exact sequence
    \[
    H_1(i_t-1)_{k_{i_t}+m_{i_t}} \longrightarrow H_1(i_t)_{k_{i_t}+m_{i_t}} \longrightarrow (0:_{M\langle i_t-1 \rangle} \ell_{i_t})_{m_{i_t}}\longrightarrow 0.
    \]
    Since the left homology group is $0$, (b)(ii) then follows for $j=i_t$.

    Now suppose $j > i_t$, and consider the exact sequence
    \begin{align*}
    &H_{j-i_t+1}(j-1)_{k_j+\cdots + k_{i_t} + m_{i_t}} \longrightarrow H_{j-i_t+1}(j)_{k_j+\cdots + k_{i_t} + m_{i_t}} \longrightarrow H_{j-i_t}(j-1)_{k_{j-1}+\cdots + k_{i_t} + m_{i_t}}\\
    &\longrightarrow H_{j-i_t}(j-1)_{k_j+\cdots + k_{i_t} + m_{i_t}}.
    \end{align*}
    It follows directly from (b)(i) that the left homology group is $0$. The right homology group is also $0$: if $j-1<i_{t+1}$, then $m_{j-1}=m_{i_t}$ and it follows from the definition of $m_{j-1}$; if $j-1\geq i_{t+1}$, it follows directly from (b)(i). Thus, $H_{j-i_t+1}(j)_{k_j+\cdots+k_{i_t}+m_{i_t}} \cong H_{j-i_t}(j-1)_{k_{j-1}+\cdots + k_{i_t} + m_{i_t}}$, which completes the inductive proof.
\end{proof}

\end{document}
